\section{Pacotes e personalização}

\begin{frame}[fragile]{Pacotes usados neste modelo}
\begin{lstlisting}
\usepackage{graphicx}  % imagens
\usepackage{booktabs}  % tabelas
\usepackage{listings}  % codigo-fonte
\usepackage[brazilian]{babel}
\usepackage{csquotes}  % citacoes
\usepackage[style=abnt]{biblatex}
\end{lstlisting}

Cada pacote acrescenta comandos ou altera regras de composição. As chamadas ficam no preâmbulo, antes de \verb|\begin{document}|.
\end{frame}

\begin{frame}[fragile]{Tema, logotipo e rodapé}
\begin{lstlisting}
\usetheme{Warsaw}
\logo{\includegraphics[height=.5cm]{images/unesp.png}}
\setbeamertemplate{headline}{}
\setbeamertemplate{frametitle continuation}{}
\end{lstlisting}

\begin{itemize}
    \item \verb|\usetheme| seleciona o tema visual do Beamer.
    \item \verb|\logo| adiciona uma imagem aos slides do tema.
    \item \verb|\setbeamertemplate| substitui um componente visual.
\end{itemize}
\end{frame}

\begin{frame}[fragile]{Espaçamento e alinhamento}
\begin{lstlisting}
\centering                 % centraliza
\raggedright               % alinha a esquerda
\vspace{1em}               % espaco vertical
\hspace{.5cm}              % espaco horizontal
\includegraphics[width=.8\linewidth]{arquivo}
\end{lstlisting}

Prefira espaçamento controlado pelo próprio ambiente. Use \verb|\vspace| e \verb|\hspace| apenas para pequenos ajustes locais.
\end{frame}
